% =====================================================================
% Frente 2 --- Caminhos e proximidade (Vitória)
% =====================================================================
\section{Frente 2 --- Grau, caminhos e proximidade}

\noindent\textbf{Responsável:} Vitória.
\textbf{No relatório:} trabalhos relacionados; caminhos e proximidade.
\textbf{Na Etapa 1:} plano de desenvolvimento com os marcos.

\subsection{O que foi implementado}

{\small\begin{tabularx}{\linewidth}{@{}L{3.9cm} L{4.7cm} Y L{1.9cm}@{}}
\toprule
Função & Arquivo & Complexidade & Situação \\
\midrule
\codigo{centralidade\_de\_grau}, \codigo{distribuicao\_de\_graus} & \codigo{metricas/grau.py} & $O(n)$ & \andamento \\
\codigo{densidade}, \codigo{grau\_medio} & \codigo{metricas/grau.py} & $O(1)$ & \andamento \\
\codigo{bfs}, \codigo{ordem\_dfs} & \codigo{metricas/caminhos.py} & $O(n+m)$ & \afazer \\
\codigo{componentes\_conexas} & \codigo{metricas/caminhos.py} & $O(n+m)$ & \afazer \\
\codigo{dijkstra} & \codigo{metricas/caminhos.py} & $O((n+m)\log n)$~\cite{dijkstra1959} & \afazer \\
\codigo{distancia\_media}, \codigo{diametro} & \codigo{metricas/caminhos.py} & $O(n(n+m))$ & \afazer \\
\codigo{proximidade}, \codigo{harmonica} & \codigo{metricas/proximidade.py} & $O(n(n+m))$~\cite{boldi2014} & \afazer \\
\bottomrule
\end{tabularx}}

\subsection{Decisões e justificativas}
\begin{itemize}[leftmargin=*]
  \item \pendente{ex.: BFS com \codigo{collections.deque}, porque \codigo{list.pop(0)} custa $O(n)$}
  \item \pendente{ex.: DFS iterativa, porque o limite de recursão do Python estoura em caminhos longos}
\end{itemize}

\subsection{Corretude: testes e resultados}

{\small\begin{tabularx}{\linewidth}{@{}Y L{2.5cm} L{2.2cm} L{3.8cm}@{}}
\toprule
Teste & Base & Obtido & Referência \\
\midrule
Densidade, grau médio e grau máximo & Karate & \falta & 0,139; 4,588; 17 \\
Distância média e diâmetro & Karate & \falta & 2,4082 e 5 \\
Closeness dos nós 0, 2 e 33 & Karate & \falta & 0,569; 0,5593; 0,55 \\
Harmônica dos nós 33 e 0 & Karate & \falta & 23,25 e 23,167 \\
Componentes, distância média e diâmetro & Twitch PT & \falta & 1; 2,5324; 7 \\
Closeness máxima & Twitch PT & \falta & nó 127 (0,6009) \\
\bottomrule
\end{tabularx}}

\subsection{Figuras e tabelas para o relatório}
\begin{itemize}[leftmargin=*]
  \item \pendente{ex.: histograma da distribuição de graus da Twitch PT}
\end{itemize}

\subsection{Limitações e pendências}
\begin{itemize}[leftmargin=*]
  \item \pendente{ex.: distância média exata inviável no ego-Facebook; usar amostra de origens}
\end{itemize}

\subsection{Registro}
\begin{description}[leftmargin=2.6cm,style=nextline]
  \registro{2026-11-11}{Métricas de grau e 13 testes (PR \#6, em revisão).}
\end{description}
